\cleardoublepage

\section{相关工作与技术背景}

本章介绍本文后续问题建模、方法设计和实验分析所依赖的背景知识与相关研究。按照本文的研究目标，本章重点围绕混合专家模型的稀疏路由机制、大语言模型推理系统中的显存管理、资源受限场景下的权重卸载，以及 MoE 专家缓存与预取优化展开。本章内容服务于后文的专家访问序列建模和 WREC 专家缓存策略设计，本文方法的具体定义与实现将在后续章节展开。

\subsection{混合专家模型与稀疏路由机制}

混合专家模型（Mixture-of-Experts, MoE）的核心思想是在神经网络中引入多个相对独立的专家子网络，并通过路由器为不同输入动态选择其中一部分专家参与计算。早期 MoE 研究已经提出由门控网络组合多个局部专家的思想~\cite{JacobsJNH91,jordan1994hierarchical}，后续稀疏门控 MoE 工作进一步表明，在只激活少量专家的条件下，可以显著扩大模型参数规模，从而提高模型容量而不线性增加每个词元的计算量~\cite{shazeer2017outrageously}。GShard 和 Switch Transformer 将稀疏专家结构扩展到大规模 Transformer 训练中，分别从分布式扩展和简化路由的角度说明 MoE 是大规模语言模型扩展的重要路线~\cite{lepikhin2020gshard,fedus2022switch}。

在典型 Transformer MoE 层中，原本的前馈网络（Feed-Forward Network, FFN）被替换为多个专家 FFN。给定输入词元表示 $x$，路由器首先计算该 词元对不同专家的选择分数，然后选择 top-$k$ 个专家参与计算。一个简化的 MoE 层输出可以表示为：

\begin{equation}
    \label{equ:moe-layer}
    y(x)=\sum_{i\in \mathcal{S}(x)} g_i(x)E_i(x),
\end{equation}


\begin{figure}[htbp]
    \centering
    \includegraphics[width=0.5\linewidth]{sparse_routing_process_in_moe_layer.png}
    \caption{典型 MoE 层的稀疏路由过程示意图}
    \label{fig:moe-sparse-routing}
\end{figure}

其中，$\mathcal{S}(x)$ 表示路由器为输入 $x$ 选出的专家集合，$E_i(\cdot)$ 表示第 $i$ 个专家网络，$g_i(x)$ 是对应的路由权重。该表达式说明，MoE 的计算路径由路由器动态决定；即使模型参数总量很大，单个词元 在一次 MoE 层计算中也只访问其中少数专家。

近年来，GLaM、Mixtral、DeepSeekMoE 等模型进一步表明，稀疏专家结构已经从训练扩展方法发展为大语言模型架构设计中的重要组成部分~\cite{glam,jiang2024mixtral,deepseekmoe2024}。以 Mixtral 为代表的开源 MoE 大模型通常采用 top-$2$ 路由，在每个 MoE 层为每个词元选择两个专家参与计算~\cite{jiang2024mixtral}；专家选择路由（Expert Choice Routing） 等工作则从专家侧选择词元，为缓解负载不均衡提供了另一类路由设计思路~\cite{zhou2022expertchoice}。这些设计带来了两个与系统部署直接相关的特征。第一，专家权重占据模型参数的主要部分，全部专家常驻 GPU 会对显存容量提出较大需求。第二，路由选择会把推理过程转化为词元、层和专家 共同决定的动态专家访问序列。如果系统不能使所有专家常驻 GPU，那么每一次专家访问都可能对应一次缓存查询，缓存未命中则会触发专家权重从 CPU 内存到 GPU 显存的传输。

MoE 路由机制还影响专家访问的局部性。对于一个请求内部的相邻词元，路由器可能反复选择相同或相近专家；对于相邻层，专家选择也可能存在统计上的转移关系；对于一个工作负载，不同请求之间又可能形成相对稳定的专家热度分布。本文后续的 WREC 策略正是围绕这些可被在线利用的访问结构展开，但本章只介绍其背景动机，具体统计定义和策略打分方式将在后续章节给出。

\subsection{大语言模型推理系统中的显存管理}

大语言模型推理通常包含预填充阶段（prefill）和解码阶段（decode）。预填充阶段一次性处理输入提示词，生成初始 KV cache；解码阶段逐词元生成输出，并不断追加新的 KV cache。二者对系统资源的压力不同：预填充阶段通常具有较大的并行计算量和较密集的词元访问，解码阶段则更容易受到自回归串行生成和请求调度影响。因此，推理系统需要同时管理模型权重、KV cache、临时激活和批处理调度。

连续批处理是提升大模型推理服务吞吐量的重要技术。Orca 将请求级调度推进到迭代级调度，使不同请求可以在生成过程中动态进入和离开批次，从而提高 GPU 利用率~\cite{yu2022orca}。在此基础上，后续系统进一步关注服务质量目标、长度预测和请求排序等调度问题。例如，学习排序式调度和基于过去-未来信息的调度方法尝试在吞吐、延迟和 SLA 之间取得平衡~\cite{fu2024efficient,gong2024pastfuture}。这些工作主要改变请求排队和 批次组织方式，而本文关注的对象更低一层：给定路由结果后，专家权重是否已经驻留在 GPU 上，以及缓存未命中造成的专家传输如何影响推理开销。

显存管理是推理系统的另一条关键主线。PagedAttention 将 KV cache 组织为分页块，减少 KV cache 连续分配带来的外部碎片，并支持更灵活的请求调度和内存复用~\cite{kwon2023pagedattention}。DeepSpeed-Inference 则从张量并行、算子融合、通信优化等角度提高大模型推理效率~\cite{aminabadi2022deepspeed}。SGLang 进一步从结构化语言模型程序和运行时调度角度提高服务效率~\cite{zheng2024sglang}。在算子层面，FlashAttention 系列通过改进注意力计算中的 I/O 访问模式降低显存读写开销，说明内存访问效率同样会影响大模型推理性能~\cite{dao2022flashattention,dao2023flashattention2}。这些系统说明，大模型推理优化不仅依赖单个算子速度，还依赖参数、KV cache 和执行计划在硬件资源之间的协调。

对于 MoE 模型，显存管理问题进一步复杂化。稠密模型的权重访问在每一层基本固定，而 MoE 模型的专家权重访问由路由器动态决定。当专家总权重超过可用 GPU 显存时，系统需要选择一部分专家常驻 GPU，其余专家在访问时按需加载。此时，专家缓存管理与 KV cache 管理共同竞争有限显存，但二者的访问模式不同：KV cache 与请求上下文长度紧密相关，专家权重则与路由事件序列紧密相关。因此，直接套用通用 KV cache 管理策略并不能完整解决 MoE 专家驻留问题。

\subsection{资源受限推理中的权重卸载}

当模型参数或运行时状态超过 GPU 显存容量时，权重卸载（offloading）成为资源受限推理的重要手段。FlexGen 面向单 GPU 环境，联合 GPU、CPU 和磁盘空间，为模型权重、KV cache 和激活设计静态调度与卸载策略，以提高大模型离线推理吞吐~\cite{sheng2023flexgen}。PowerInfer 观察到部分神经元在本地推理中更频繁激活，并尝试将热权重保留在 GPU 上、冷权重放置在 CPU 上，从而改善消费级 GPU 上的推理效率~\cite{song2023powerinfer}。LLM in a Flash 则关注有限内存设备上的权重加载与稀疏访问问题，通过减少数据搬移支持更大模型在内存受限设备上运行~\cite{alizadeh2024llminflash}。

这些工作与本文共享一个基本前提：在资源受限环境下，数据搬移本身会成为推理性能的重要瓶颈。不同之处在于，通用稠密模型卸载通常面对的是层级或张量级的相对固定访问顺序，而 MoE 专家卸载面对的是路由器产生的动态专家访问序列。对于稠密 Transformer，若给定模型结构和执行阶段，下一层权重访问基本可以由计算图直接确定；对于 MoE，下一次访问哪个专家取决于当前词元的路由结果，且不同请求、不同层之间可能存在专家热度差异、短期复用和跨层转移等统计关系。

资源受限推理中的另一个关键问题是传输与计算能否重叠。流水线并行、推理阶段调度以及推测解码等方向均体现了通过重叠或减少串行等待提升推理效率的思路~\cite{huang2019gpipe,narayanan2019pipedream,leviathan2023speculative}。如果权重传输可以完全隐藏在计算之后，则缓存未命中的性能影响较小；如果 CPU-GPU 带宽不足，或专家权重较大而无法及时预取，则未命中会直接转化为额外等待时间。本文后续参考 HRM 的分层瓶颈分析思想判断目标配置是否处于专家传输受限状态（专家传输受限（expert-transfer-bound）），其目的正是区分“专家传输是主要瓶颈”的场景与“专家传输可以被其他计算掩盖”的场景。在前者中，专家缓存调度的性能收益更容易体现。

\subsection{MoE 专家卸载、缓存与预取相关研究}

近年来，已有研究开始针对 MoE 模型中的专家按需加载、缓存和预取进行优化。Fast MoE offloading 相关工作围绕 Mixtral 等模型讨论了专家按需加载、缓存和异步预取问题，说明在单卡或有限显存环境下，专家权重不能简单视为始终常驻的静态参数~\cite{eliseev2023offloading}。MoE-Infinity 进一步从激活感知角度设计专家卸载策略，希望根据专家实际激活情况减少不必要的数据搬移~\cite{xue2024moeinfinity}。Fiddler 从 CPU-GPU 协同执行角度处理资源受限 MoE 推理问题，根据输入规模和批次特征在 CPU 计算专家与 GPU 加载专家之间进行选择，说明专家权重搬移并非唯一可行路径，系统也可以通过异构计算分担部分专家执行开销~\cite{kamahori2024fiddler}。MoE-Lightning 则关注内存受限 GPU 上的高吞吐 MoE 推理，强调专家放置、传输和执行调度之间的联合优化，并提出 HRM 从分层 Roofline 角度分析 CPU/GPU 计算能力、内存带宽和跨设备传输带宽之间的瓶颈关系~\cite{cao2025moelightning}。此外，面向 MoE 的通信优化、算子感知卸载和分离式推理等工作也说明，专家权重放置、跨设备传输和运行时调度已经成为 MoE 推理系统优化中的重要问题~\cite{mei2025moecomm,zhang2025operatorawareoffload,sun2026pddisaggmoe}。

上述研究说明，MoE 专家管理包含三个相互关联的问题。第一是专家驻留集合选择，即在有限 GPU 显存中保留哪些专家。第二是换出策略，即当新专家需要加载而缓存已满时，应该换出哪个专家。第三是预取策略，即能否在专家真正被访问前提前加载，减少后续等待时间。这三个问题与经典缓存问题类似，但 MoE 场景具有更强的结构性：专家访问不仅有时间顺序，还包含层号、请求编号、词元位置和 top-$k$ 路由结果等上下文。

现有 MoE 专家卸载工作通常依赖某种形式的专家热度、路由预测或执行流水线。静态热度可以利用训练集或历史工作负载统计，但当请求内部局部性较强时，单纯热度排序可能无法充分利用短期复用。最近最少使用（Least Recently Used, LRU）等在线缓存策略可以利用近期访问，但它不区分同一请求内部的结构性复用、跨层转移关系和工作负载级先验。预取策略如果能够准确预测未来专家访问，则可以减少等待时间；但如果预测错误或预取过多，也会带来额外传输和缓存污染。近期也有工作尝试从模型内部表示中预测后续专家选择，例如在注意力计算前引入轻量专家预测器以提前发起专家预取~\cite{zhu2025preattention}。与这类预测型预取不同，本文不训练额外专家预测器，而是仅利用已经观测到的路由事件和历史工作负载统计构造在线缓存评分。

也有研究从模型是否适合专家卸载的角度分析 MoE 路由一致性。相关工作指出，不同 MoE 模型的本地路由稳定性并不相同，因此并非所有 MoE 都同样适合专家卸载优化~\cite{liang2026notallmodels}。这一观察与本文的问题设定一致：在设计专家缓存策略前，需要先分析目标模型和工作负载的专家访问局部性，而不能默认任何 MoE 都存在同等可利用的缓存收益。

\subsection{缓存策略与专家访问局部性}

从抽象角度看，专家缓存可以被视为带有传输成本的在线缓存问题。系统维护一个容量有限的专家驻留集合；当访问的专家已经驻留在 GPU 上时，访问命中；当专家不在该集合中时，系统需要加载该专家，并在容量不足时选择已有专家换出。若每个专家权重大小近似相同，则缓存容量可以用可容纳的专家数量表示；若专家大小或传输代价不同，则需要进一步考虑加权容量和异构传输成本。

经典 LRU 策略假设最近访问过的对象更可能再次访问，因此在缓存满时换出最久未访问对象。静态热度策略则假设长期高频对象更值得常驻。类似地，H2O 等 KV cache 管理工作也利用访问重要性或历史统计在有限缓存中保留更有价值的对象~\cite{zhang2024h2o}。二者分别对应时间局部性和频率局部性，但 MoE 专家访问还包含请求内局部性与跨层结构。举例来说，同一请求的一段输入词元可能反复激活少数专家；某一层的专家选择也可能对下一层专家选择提供统计线索。若缓存策略只利用全局频率或最近访问时间，就可能遗漏这些上下文信息。

本文后续采用基于访问轨迹的重放分析对不同策略进行比较，其基本思想是将真实或离线采集到的路由结果整理为专家引用序列，然后在统一缓存预算下重放该序列，统计命中、未命中、传输量和传输时间代理指标。该方法有两个优点。第一，它能够在不改变模型输出的前提下隔离专家缓存策略本身的影响。第二，它允许在同一条访问序列上比较在线策略、静态策略和使用未来访问信息的离线上界策略，从而评估在线策略与离线上界之间的差距。

需要强调的是，离线最优或使用未来信息的策略只能作为上界和压力测试，不能直接视为可部署方案。Belady 离线上界源于经典缓存替换问题，在已知完整未来访问序列时换出下一次访问最晚的对象，可作为等大小缓存对象假设下的离线最优参考~\cite{belady1966study}。实际推理服务系统中，缓存策略只能使用当前及过去可见的信息，或者使用来自训练集、历史工作负载的统计先验。因此，本文在相关实验中区分公平在线基线、静态先验基线和离线上界参考结果，并将 WREC 定位为不改变路由结果、仅使用可在线维护状态的专家缓存调度方法。

\subsection{本文工作的定位}

综合以上研究现状，资源受限 MoE 推理优化可以大致分为表~\ref{tab:related-work-position} 所示的几类。通用 LLM 推理服务系统主要优化批处理调度、KV cache 和执行流水线；资源受限卸载系统关注在 GPU、CPU 和磁盘之间安排权重与状态；MoE 专家卸载系统进一步关注专家级驻留集合、换出策略和预取策略；而本文的工作聚焦于专家传输受限、部分专家驻留场景下的预填充阶段专家缓存调度。

\begin{table}[ht]
    \centering
    \caption{\label{tab:related-work-position}本文与相关研究方向的关系}
    \begin{tabularx}{\linewidth}{lXX}
        \hline
        研究类别 & 代表性问题 & 与本文研究的关系 \\ \hline
        通用 LLM 推理服务 & 连续批处理、KV cache 管理、请求调度与执行流水线~\cite{yu2022orca,kwon2023pagedattention} & 主要关注请求调度和 KV cache 管理，未直接面向专家驻留集合选择问题。 \\ \hline
        资源受限卸载 & 在 GPU、CPU 或磁盘之间安排权重、KV cache 和激活~\cite{sheng2023flexgen,song2023powerinfer} & 说明数据搬移可能成为性能瓶颈，为本文分析专家传输开销提供依据。 \\ \hline
        MoE 专家卸载 & 专家按需加载、异步预取、专家放置和高吞吐执行~\cite{eliseev2023offloading,xue2024moeinfinity,kamahori2024fiddler,cao2025moelightning} & 与本文研究对象最接近；本文侧重在真实路由轨迹下比较在线缓存策略与离线上界。 \\ \hline
        路由局部性分析 & 分析不同 MoE 模型是否具备可利用的局部路由一致性~\cite{liang2026notallmodels} & 支持本文先分析路由局部性、再设计缓存策略的研究思路。 \\ \hline
    \end{tabularx}
\end{table}

与直接优化模型结构、修改路由器或重新训练 MoE 不同，本文将模型路由结果视为给定输入，研究在不改变模型输出的前提下如何降低专家权重搬移开销。与只报告系统端到端吞吐的工作不同，本文首先使用基于访问轨迹的重放模拟器隔离分析专家缓存策略本身的贡献，并参考已有 HRM 思想限定结论适用的专家传输受限场景。基于上述定位，后续章节将围绕三个问题展开：目标模型与硬件配置是否存在专家传输瓶颈，路由轨迹中是否包含可利用的访问局部性，以及 WREC 在线策略相对于公平在线基线能否进一步接近离线上界。
